\begin{document}
\docTitle{Statistical tests}
\subsection*{Paired permutation test}
The theory behind paired permutation test is based on exchangeability under the null hypothesis.
%
Suppose we have n paired observations, where each pair consists of two measurements taken under two different conditions, X and Y, while all other relevant factors are kept the same.
\[
(X_1, Y_1), (X_2, Y_2), \ldots, (X_n, Y_n)
\]
For each pair, compute the difference:

\[
D_i = Y_i - X_i
\]

The null hypothesis is usually:

\[
H_0: \text{there is no systematic difference between } X \text{ and } Y
\]

In a paired setting, “no systematic difference” means that within each pair, the labels \(X_i\) and \(Y_i\) are arbitrary. In other words, under the null, it should not matter which value is called “before” and which is called “after.”

%
%
So for each pair, the observed difference \(D_i\) could just as plausibly have been \(-D_i\). That is the key symmetry assumption:

\[
D_i \overset{d}{=} -D_i \quad \text{under } H_0
\]

So the test generates all possible sign flips:

\[
(D_1, D_2, \ldots, D_n)
\]

\[
(\pm D_1, \pm D_2, \ldots, \pm D_n)
\]

There are \(2^n\) possible sign-flipped datasets. For each one, you compute a test statistic, often the mean difference:

\[
T = \frac{1}{n}\sum_{i=1}^n D_i
\]

This gives the exact null distribution of the statistic under the assumption that the signs are arbitrary.
Let \(T_{\text{obs}}\) be the observed statistic, i.e \[T_{\text{obs}} = \frac{1}{n}\sum_{i=1}^n D_i.\]
In order to perform this test, you would compare \(T_{\text{obs}}\) to the null distribution.
The null distribution can be created by considering all possible sign flips of the observed differences. 
Specifically, for each sign-flipped dataset, \((D_1^*,D_2^*,\ldots,D_n^*)=(s_1D_1,s_2D_2,\ldots,s_nD_n)\), where each \(s_i\) is either \(+1\) or \(-1\), compute the corresponding test statistic \(T_{\text{perm}} = \frac{1}{n}\sum_{i=1}^n D_i^*\).
The collection of all such values of \(T_{\text{perm}}\) forms the permutation null distribution.
This distribution represents the values of the test statistic that could have occurred if the null hypothesis were true and the signs of the paired differences were arbitrary.
Then the observed statistic \(T_{\text{obs}}\) is compared to this null distribution. For a two-sided test, the p-value is
\[
p = \Pr\left(|T_{\text{perm}}| \ge |T_{\text{obs}}| \mid H_0\right).
\]

A small p-value means that the observed mean difference is unusual under the null hypothesis. 
Therefore, there is evidence against \(H_0\), suggesting a systematic difference between \(X\) and \(Y\).
\agentProposedEdit[created=2026-10-15T15:41:49.029Z,conversation_id=308253e6-d399-5c89-b3dd-4b6958845ab1]{agent-bda44393-eb3b-4dae-bae1-2d20f06c48ba}{Add the exact McNemar test (as the paired permutation test on binary outcomes), the Holm correction and the paired bootstrap, the three procedures used in the data-preprocessing paper's BasicSkills tables.
}{}{
\subsection*{Exact McNemar test}
The exact McNemar test is the paired permutation test above, applied to binary outcomes.
Suppose two models A and B answer the same \(n\) questions, and let \(X_i, Y_i \in \{0,1\}\) be 1 when A, respectively B, answers question \(i\) correctly.
The paired differences then take only three values:
\[
D_i = Y_i - X_i \in \{-1, 0, +1\}.
\]
The \(n\) questions fall into four counts:
\[
\begin{array}{l|cc}
 & \text{B correct} & \text{B wrong} \\
\hline
\text{A correct} & a & b \\
\text{A wrong} & c & d
\end{array}
\]
The mean difference is exactly the difference in accuracy:
\[
T_{\text{obs}} = \frac{1}{n}\sum_{i=1}^n D_i = \frac{c-b}{n} = \text{acc}_B - \text{acc}_A.
\]
The \(a+d\) concordant questions, which both models get right or both get wrong, have \(D_i = 0\), and flipping their sign changes nothing.
So they carry no information about which model is better, and only the \(n_d = b + c\) discordant questions matter.
Under \(H_0\), the sign of each discordant difference is a fair coin, so the number \(C\) of discordant questions won by B follows
\[
C \sim \text{Binomial}\left(b+c, \tfrac12\right) \quad \text{under } H_0.
\]
This binomial is the permutation null distribution of the previous section in closed form, so all \(2^n\) sign flips are covered without enumerating them.
Since \(|T_{\text{perm}}| \ge |T_{\text{obs}}|\) exactly when \(C \le \min(b,c)\) or \(C \ge \max(b,c)\), the two-sided p-value is
\[
p = \min\left(1,\; 2\sum_{k=0}^{\min(b,c)} \binom{b+c}{k} 2^{-(b+c)}\right).
\]
For example, with \(b = 3\) and \(c = 12\), \(\Pr(C \le 3) = (1 + 15 + 105 + 455)/2^{15} \approx 0.018\), so \(p \approx 0.035\).
The classical McNemar test replaces the binomial by its large-sample approximation \(\chi^2 = (b-c)^2/(b+c)\) with one degree of freedom; the exact version is preferred when \(b + c\) is small.

Pairing is what makes the test sensitive.
Comparing the two accuracies as if they came from independent samples would treat the variation in question difficulty as noise, although both models face the same questions.
In the paired test, a question that is easy or hard for both models lands in \(a\) or \(d\) and drops out, so only the questions on which the models actually disagree decide the result.

\subsection*{Holm correction for multiple tests}
When \(m\) hypotheses are each tested at level \(\alpha\), the probability of at least one false rejection, the family-wise error rate (FWER), can be much larger than \(\alpha\).
With \(m = 12\) independent true null hypotheses and \(\alpha = 0.05\), it is \(1 - 0.95^{12} \approx 0.46\).
The Bonferroni correction rejects \(H_j\) only when \(p_j \le \alpha/m\), which keeps the FWER at most \(\alpha\) by the union bound.
The Holm correction is a step-down version that rejects at least as much and needs no assumption about how the tests depend on each other.
Sort the p-values as \(p_{(1)} \le p_{(2)} \le \cdots \le p_{(m)}\) and compare them with increasingly lenient thresholds:
\[
\text{for } k = 1, 2, \ldots, m: \quad \text{reject } H_{(k)} \text{ if } p_{(k)} \le \frac{\alpha}{m-k+1}, \text{ otherwise stop and keep } H_{(k)}, \ldots, H_{(m)}.
\]
Tables usually report the equivalent Holm-adjusted p-values, which are compared directly with \(\alpha\):
\[
\tilde p_{(k)} = \min\left(1,\; \max_{j \le k}\, (m-j+1)\, p_{(j)}\right).
\]
Why the FWER stays at most \(\alpha\): let \(m_0\) of the hypotheses be true.
By the union bound, with probability at least \(1-\alpha\) every true null has \(p > \alpha/m_0\).
At most \(m - m_0\) false nulls can be ranked before the smallest true-null p-value, so its rank is \(k \le m - m_0 + 1\) and its threshold is \(\alpha/(m-k+1) \le \alpha/m_0\).
On that event it fails its threshold, the procedure stops there, and no true null is rejected.

\subsection*{Paired bootstrap}
Some statistics are not a single comparison per question, for example the mean of per-skill accuracies over several skills, where each skill has its own number of questions.
For such statistics, the paired bootstrap estimates how much the statistic would vary if the questions had been drawn again.
Let \(\theta\) be the statistic of interest, here the difference between model B's and model A's mean per-skill accuracy, and \(\hat\theta\) its observed value.
For \(r = 1, \ldots, R\) (typically \(R = 10{,}000\)), draw questions with replacement, within each skill so that every skill keeps its size, keep each question's pair \((X_i, Y_i)\) together, and recompute the statistic \(\hat\theta^*_r\) on the resample.
The 95\% percentile confidence interval is the range between the 2.5\% and 97.5\% quantiles of \(\hat\theta^*_1, \ldots, \hat\theta^*_R\).
A common two-sided p-value for \(H_0: \theta = 0\) is
\[
p = \min\left(1,\; 2\min\left(\Pr{}^*(\hat\theta^* \le 0),\; \Pr{}^*(\hat\theta^* \ge 0)\right)\right),
\]
where \(\Pr^*\) is the fraction of resamples; it is the smallest \(\alpha\) for which the \((1-\alpha)\) percentile interval excludes 0.
Keeping the pairs together preserves the correlation between the two models' answers, so shared question difficulty cancels as in the McNemar test; resampling the two models independently would give a wider interval.
Unlike the permutation and McNemar tests, which build the distribution of the statistic under \(H_0\), the bootstrap approximates its sampling distribution around the observed value, so its p-values are approximate rather than exact.
}\end{document}
